%==========================================================
% AUTHOR'S NOTE — Volume I: Zero to Bit
%==========================================================
{\small \linespread{1.35}\selectfont
\chapter*{A Note on This Volume}
\addcontentsline{toc}{chapter}{A Note on This Volume}
\chaptermark{A Note on This Volume}

This volume is a first pass, not a complete treatment. The chapters on
integer representation and floating-point arithmetic are thorough because
everything else in this work depends on them. The chapters on serialization
and compression are shorter --- enough to know what the topic is and what
questions it raises, not enough to replace a dedicated reference. That
imbalance is deliberate.

A few things worth knowing before you continue.

\medskip

\noindent\textbf{On the code.}\quad
Where this volume includes code examples, they are written to be read, not
to be fast. If a chapter discusses an operation that has a well-known
library implementation, the example still builds it from scratch. The
purpose is understanding, not performance. When you want the tuned version,
the relevant header is in Appendix~A.

\medskip

\noindent\textbf{On assumptions.}\quad
This volume assumes \texttt{CHAR\_BIT == 8} and IEEE~754 binary
floating-point throughout. Both hold on every common general-purpose
processor. If you are working on hardware that violates these assumptions,
the concepts still apply; the specific bit patterns and boundary values
may differ.

\medskip

\noindent\textbf{On scope.}\quad
Some topics that naturally belong to this layer of the stack --- Unicode
normalization algorithms, arbitrary-precision arithmetic, advanced entropy
coding --- are introduced but not fully developed. This is not an oversight.
Those subjects are large enough to occupy their own books, and covering
them completely here would shift the focus away from arrays. The citations
at the end point to the right sources for each one.

\medskip

\noindent\textbf{On errors.}\quad
Some mistakes remain in this text. This is a living document and it
improves over time, but it does not pretend to be finished. If you find
something wrong --- a calculation, an explanation, a claim about a
standard --- please open an issue at
\url{https://github.com/papyrxis/arliz}. Corrections are read and
applied.\\
\medskip
Volume~I is preparatory. None of it is self-contained in the sense that
understanding representation is a complete goal; it becomes meaningful when
you can see what the hardware does with it, and what arrays do with both.
That is what the next two volumes are for.

\begin{flushright}
  \textit{Mahdi Mamashli}\\
  \textit{Iran, 2026}
\end{flushright}

\clearpage
}